\documentclass[11pt,a4paper]{article}
\usepackage[T1]{fontenc}
\renewcommand{\rmdefault}{lmr}
\renewcommand{\ttdefault}{lmtt}
\usepackage[margin=27mm]{geometry}
\usepackage{amsmath,amssymb,amsthm}
\usepackage{array,longtable}
\usepackage[hidelinks]{hyperref}
\usepackage{url}
\setlength{\emergencystretch}{2em}
\numberwithin{equation}{section}
\newtheorem{theorem}{Theorem}[section]
\newtheorem{proposition}[theorem]{Proposition}
\newtheorem{lemma}[theorem]{Lemma}
\newtheorem{corollary}[theorem]{Corollary}
\theoremstyle{remark}
\newtheorem{remark}[theorem]{Remark}
\DeclareMathOperator{\Tr}{Tr}
\DeclareMathOperator{\diag}{diag}
\DeclareMathOperator{\spanof}{span}
\DeclareMathOperator{\Var}{Var}
\DeclareMathOperator{\CV}{CV}
\newcommand{\CK}{\mathcal C_K}
\newcommand{\CS}{\mathcal C_S}
\newcommand{\CM}{C_{\mathrm M}}
\newcommand{\SU}{S_{U^*}}
\newcommand{\SI}{S_I}
\newcommand{\KI}{K_I}
\newcommand{\KU}{K_{U^*}}
\newcommand{\HS}{\mathrm{HS}}
\newcommand{\vecr}{\operatorname{vec}_{\mathrm r}}
\newcommand{\eps}{\epsilon}
\newcommand{\avg}[1]{\left\langle #1\right\rangle_T}
\newcommand{\ket}[1]{\lvert #1\rangle}
\newcommand{\bra}[1]{\langle #1\rvert}
\newcommand{\norm}[1]{\left\lVert #1\right\rVert}
\title{Exact limits of Krylov-complexity comparisons under purification}
% Human author; no ORCID supplied.
\author{Dehao Lin\\{\small School of Physics, Sun Yat-sen University, Guangzhou, China}}
\date{}
\hypersetup{pdftitle={Exact limits of Krylov-complexity comparisons under purification},pdfauthor={Dehao Lin},pdfsubject={Final public v0.1 package; not yet published; authoritative V3 scientific freeze}}
\begin{document}
\maketitle
\begin{abstract}
Purification relates mixed quantum states to pure states on a larger Hilbert
space, but it changes the spectral weights that determine Krylov complexity.
We examine the comparison conjectures of Das and Mori using their normalized
Hilbert--Schmidt and spread-complexity conventions. Each side of their
two-sided hierarchy admits an exact full-rank qutrit counterexample, with
separate rational finite-time certificates. The hierarchy holds for every
qubit, so physical dimension three is minimal for failure of either side.
At fixed dimension four and fixed Hamiltonian, two different full-rank
families have unbounded recurrence-period coefficients of variation for the
purification-to-mixed complexity ratio and its reciprocal, respectively.
A third family has fixed purity and an unbounded ratio and temporal mean.
These are no-go theorems for explicit universal concentration and purity-only
control formulations motivated by the qualitative second conjecture.
In contrast, the pure-density factor-two inequality holds in finite dimension;
we give a self-contained proof and identify it as a corollary of existing
tensor-product superadditivity. A qubit concentration bound and
state-dependent return bounds delimit the surviving comparisons. The common
mechanism is the distinct weighting of coherent sectors by the density matrix
and its positive square root.
\end{abstract}
\section{Introduction and source problem}\label{sec:introduction}
Krylov constructions convert unitary evolution into motion along a basis
generated by successive applications of the dynamical generator. Operator
complexity measures the mean position in an operator Krylov chain, while
spread complexity uses a state-vector chain
\cite{Parker2019,Balasubramanian2022}. A density matrix can participate in
either construction: it is itself an operator, and it also specifies a pure
state after purification. Comparing these descriptions requires control of
both the seed normalization and the generator. Earlier work on density-matrix
operators already distinguishes their Krylov complexity from the spread
complexity of a pure state \cite{Caputa2024}.

Das and Mori \cite{DasMori2026} formulated this comparison for two unitary
purification schemes. The time-independent scheme evolves only the physical
system, while the time-dependent scheme evolves the ancilla by the conjugate
unitary. Their Eq.~(4) proposes
\begin{equation}\label{eq:source-hierarchy}
 \CS\bigl(\Psi_\rho^{U^*}(t)\bigr)
 \ \le\ \CK\bigl(\rho(t)\bigr)
 \ \le\ \CK\bigl(\Psi_\rho^I(t)\bigr).
\end{equation}
Here the rightmost quantity is the operator complexity of the rank-one
density matrix of the purification. It is not the spread complexity of the
purified vector. We retain the published source's conventions throughout.

The source also proposes two different kinds of statement. Its second
conjecture describes the ratio of the time-dependent purification's spread
complexity to the mixed-state operator complexity as approximately constant
relative to its temporal mean, with a value depending on purity. This is
qualitative: it specifies neither a uniform tolerance nor an exact function
of purity. The third conjecture is the precise factor-two inequality between
operator and spread complexities of the same purification. The distinction
between these statements is central to the analysis below.

We give separate exact qutrit witnesses against the two sides of
Eq.~\eqref{eq:source-hierarchy}. The source's qubit hierarchy remains valid;
its explicit three-atom reduction proves that neither inequality can fail in
a smaller physical dimension. For temporal ratios, we state and rule out
three universal formulations: a finite uniform recurrence-period coefficient
of variation for the main ratio, the corresponding bound for its reciprocal,
and finite purity-only upper control of the main ratio or its mean. The
two concentration questions are logically distinct, so they require different
families. All counterexamples use strictly positive density matrices at every
allowed parameter value.

The factor-two inequality has a different status. Tensor-product
superadditivity, stated by Murugan and van Zyl
\cite[Sec.~2, Eq.~(2.4)]{MuruganVanZyl2026}, implies it after vectorizing a
rank-one projector. We give the needed finite-dimensional proof in terms of
nested polynomial subspaces and apply it separately to the two purification
generators. This known implication is context for the failed mixed-state
comparisons, rather than a new headline inequality. The early-time
pure-density relation also has antecedents in Ref.~\cite{Caputa2024}.

Our focus is finite-dimensional, time-independent unitary dynamics. The
results do not assess nonunitary examples or typical-state limits. A final fresh
public-prior-art sweep before release found no matching
public prior result; absolute priority
is not claimed.

Section~\ref{sec:framework} fixes the definitions and spectral representation.
Sections~\ref{sec:hierarchy}--\ref{sec:ratios} contain the comparison results
and their proofs. Section~\ref{sec:survivors} gives restricted survivors and
the spectral mechanism. Exact finite certificates and technical estimates
are collected in the appendices.

\section{Definitions and spectral measures}\label{sec:framework}
\subsection{Seeds, generators, and normalization}
Let $H$ be a time-independent Hermitian matrix on a physical Hilbert space
of dimension $d$, and let $\rho$ be positive semidefinite with trace one.
We set $\hbar=1$ and use
\begin{equation}\label{eq:evolution}
 U_t=e^{-itH},\qquad \rho(t)=U_t\rho U_t^\dagger,
 \qquad P=\Tr\rho^2.
\end{equation}
For a Hermitian generator $G$ and a unit vector $v$, orthonormalize
$v,Gv,G^2v,\ldots$ until exact termination. If $K_n$ is this basis, the
complexity associated with the pair $(G,v)$ is
\begin{equation}\label{eq:general-complexity}
 C_{G,v}(t)=\sum_{n=0}^{m-1}n\,|\varphi_n(t)|^2,
 \qquad \varphi_n(t)=\bra{K_n}e^{-itG}\ket{v},
 \qquad \sum_{n=0}^{m-1}|\varphi_n(t)|^2=1.
\end{equation}
The integer $m$ is the cyclic dimension, not the physical dimension. The
Lanczos phase convention does not affect these probabilities.

For operator complexity the inner product is
$(A|B)=\Tr(A^\dagger B)$, and the generator is the self-adjoint
commutator map $\mathcal L_H A=[H,A]$. The normalized mixed seed is
$\rho/\sqrt P$. In particular, unit trace does not imply unit
Hilbert--Schmidt norm. This normalization agrees with
Ref.~\cite[Supplemental Eqs.~(S.7)--(S.11)]{DasMori2026}.

For full-rank $\rho$, fix row vectorization and define
\begin{equation}\label{eq:purification}
 \begin{aligned}
 \vecr A&=\sum_{a,b}A_{ab}\ket a\otimes\ket b,
 &s&=\vecr\sqrt\rho, &\norm{s}^2&=\Tr\rho=1,\\
 G_I&=H\otimes I,
 &G_{U^*}&=H\otimes I-I\otimes\overline H,
 &s_X(t)&=e^{-itG_X}s .
 \end{aligned}
\end{equation}
Thus $s_I(t)=\Psi_\rho^I(t)$ and
$s_{U^*}(t)=\vecr\sqrt{\rho(t)}$. A bar denotes entrywise conjugation;
for Hermitian $H$, $\overline H=H^T$. Both purifications have physical
reduced density matrix $\rho(t)$. We abbreviate
\begin{equation}\label{eq:abbreviations}
 \CM=\CK(\rho(t)),\qquad
 S_X=\CS(s_X(t)),\qquad
 K_X=\CK\bigl(\ket{s_X(t)}\bra{s_X(t)}\bigr),
 \quad X\in\{I,U^*\}.
\end{equation}
The seed for $K_X$ is $ss^\dagger$, with norm one, and its generator is
$[G_X,\cdot]$. The source hierarchy is $\SU\le\CM\le\KI$.
For rank-deficient states, statements about arbitrary purifications below
use the actual normalized pure seed and its generator on a fixed enlarged
space. No identification of canonical and minimal purifications is assumed.

\subsection{The spectral representation}
The pair $(G,v)$ determines a probability measure $\nu$ on the eigenvalues
of $G$. Orthonormal polynomials $p_n$ with respect to $\nu$ represent the
Krylov vectors, and
\begin{equation}\label{eq:spectral-formula}
 \varphi_n(t)=\int p_n(x)e^{-itx}\,d\nu(x).
\end{equation}
One may choose the polynomials real. Repeated spectral values are merged;
the number of distinct atoms with positive weight is the exact chain length.
In an energy basis of $H$, the mixed and time-dependent spread measures are
\begin{equation}\label{eq:gap-measures}
 \nu_{\mathrm M}=\sum_{a,b}\frac{|\rho_{ab}|^2}{P}\,
       \delta_{E_a-E_b},\qquad
 \nu_{S,U^*}=\sum_{a,b}|(\sqrt\rho)_{ab}|^2\,
       \delta_{E_a-E_b}.
\end{equation}
For the time-independent purified vector the energy measure is
$\nu_{S,I}=\sum_a\rho_{aa}\delta_{E_a}$. Its rank-one operator measure
is the difference convolution of this measure with itself:
\begin{equation}\label{eq:difference-convolution}
 \nu_{K,I}=\sum_{a,b}\rho_{aa}\rho_{bb}\delta_{E_a-E_b}.
\end{equation}
More generally, a pure-state measure produces the difference convolution
when its rank-one density matrix is vectorized. In contrast, the two measures
in Eq.~\eqref{eq:gap-measures} have the same possible frequencies but
different weights. Equality of frequency support does not order their
complexities.

\begin{lemma}[Three-atom chain]\label{lem:three}
For $0<\mu<1$, the measure
$\nu_\mu=(1-\mu)\delta_0+(\mu/2)(\delta_{-1}+\delta_1)$ has
complexity
\begin{equation}\label{eq:three-function}
 F(\mu;t)=\mu\sin^2t+8\mu(1-\mu)\sin^4(t/2).
\end{equation}
The same expression holds at $\mu=0,1$ with the corresponding terminated
chains. A frequency gap $\omega$ replaces $t$ by $\omega t$.
\end{lemma}
\begin{proof}
The squared nonzero Lanczos coefficients are $b_1^2=\mu$ and
$b_2^2=1-\mu$, and all diagonal coefficients vanish. The Jacobi matrix
satisfies $J^3=J$. Its evolved seed has amplitudes
\begin{equation}\label{eq:three-amplitudes}
 \bigl(1-\mu+\mu\cos t,\;-i\sqrt\mu\sin t,\;
 \sqrt{\mu(1-\mu)}(\cos t-1)\bigr).
\end{equation}
Their squared moduli sum to one, and their index-weighted sum gives
Eq.~\eqref{eq:three-function}. At either endpoint the zero-coupling branch
is simply omitted.
\end{proof}

\begin{lemma}[Return bound]\label{lem:return}
For a normalized finite chain of length $m$ with return amplitude $A(t)$,
\begin{equation}\label{eq:return-bound}
 1-|A(t)|^2\le C(t)\le(m-1)\bigl(1-|A(t)|^2\bigr).
\end{equation}
For the positive-matrix flows of $\rho$ or $\sqrt\rho$, $A(t)$ is real
and lies in $[0,1]$. If $m\le5$ and $D(t)=1-A(t)$, then
\begin{equation}\label{eq:loss-bound}
 D(t)\le C(t)\le8D(t).
\end{equation}
\end{lemma}
\begin{proof}
The total probability away from node zero is $1-|A|^2$, and each of its
indices lies between $1$ and $m-1$. Positivity of the matrix overlap and
Cauchy--Schwarz give $0\le A\le1$ in the stated flows. Finally,
$1-A^2=(1-A)(1+A)$ lies between $D$ and $2D$.
\end{proof}

\section{Exact failures of the two-sided hierarchy}\label{sec:hierarchy}
\subsection{The lower inequality}
\begin{proposition}[Full-rank qutrit lower-side witness]\label{prop:left}
For
\begin{equation}\label{eq:left-state}
 H_L=\diag(2,0,1),\qquad
 \rho_L=\frac1{42}\begin{pmatrix}32&0&0\\0&5&3\\0&3&5\end{pmatrix},
\end{equation}
the lower inequality in Eq.~\eqref{eq:source-hierarchy} fails at $t=\pi$:
\begin{equation}\label{eq:left-values}
 \SU(\pi)=\frac{82}{441},\qquad
 \CM(\pi)=\frac{1074}{8281},\qquad
 \SU(\pi)-\CM(\pi)=\frac{4192}{74529}>0.
\end{equation}
\end{proposition}
\begin{proof}
The eigenvalues of $\rho_L$ are $16/21,4/21,1/21$, all positive, and
\begin{equation}\label{eq:left-root}
 \sqrt{\rho_L}=\frac1{2\sqrt{21}}
 \begin{pmatrix}8&0&0\\0&3&1\\0&1&3\end{pmatrix},
 \qquad P_L=\frac{13}{21}.
\end{equation}
Only the second and third levels carry coherence, with energy gap one.
Equation~\eqref{eq:gap-measures} therefore reduces both problems to
Lemma~\ref{lem:three}, with
\begin{equation}\label{eq:left-weights}
 \mu_S=\frac1{42},\qquad \mu_K=\frac3{182},\qquad
 \SU(t)=F(\mu_S;t),\qquad \CM(t)=F(\mu_K;t).
\end{equation}
Evaluation at $\pi$ gives Eq.~\eqref{eq:left-values} exactly. The full
probability vectors appear in Appendix~\ref{app:certificates}.
\end{proof}

The witness involves neither a rank-deficient limit nor a numerical
truncation. The light coherent sector carries a larger relative weight in
$\sqrt\rho$ than in $\rho/\sqrt P$. Appendix~\ref{app:unbounded-left}
embeds it in a full-rank qutrit family for which $\SU(\pi)/\CM(\pi)$ is
unbounded, ruling out a repair of the lower side by a fixed positive
state-independent multiplier.

\subsection{The upper inequality}
\begin{proposition}[Full-rank qutrit upper-side witness]\label{prop:right}
For
\begin{equation}\label{eq:right-state}
 \rho_R=\frac1{26}\diag(16,1,9),\qquad
 H_R=\begin{pmatrix}0&1&0\\1&0&0\\0&0&0\end{pmatrix},
\end{equation}
the upper inequality in Eq.~\eqref{eq:source-hierarchy} fails at $t=\pi/3$:
\begin{equation}\label{eq:right-values}
 \CM(\pi/3)=\frac{1141425}{913952},\qquad
 \KI(\pi/3)=\frac{2967537}{2456246},\qquad
 \CM(\pi/3)-\KI(\pi/3)=\frac{1600683}{39299936}>0.
\end{equation}
It also fails throughout the explicit interval $0<t<1/2704$.
\end{proposition}
\begin{proof}
The state is strictly positive, has purity $P_R=1/2$, and has positive
square root $\diag(4,1,3)/\sqrt{26}$. In the energy basis of $H_R$, the
normalized mixed-operator measure and the purified-operator measure are
\begin{equation}\label{eq:right-measures}
 \begin{array}{c|c|c}
 &\text{ordered spectral support}&\text{ordered weights}\\ \hline
 \CM&(-2,0,2)&(225/1352,\ 451/676,\ 225/1352)\\[2pt]
 \KI&(-2,-1,0,1,2)&(289/2704,\ 153/676,\ 451/1352,\ 153/676,\ 289/2704)
 \end{array}
\end{equation}
For the second row the purified vector has energy weights $17/52,9/26,17/52$
on $-1,0,1$; taking their difference convolution gives the stated row.
All weights are positive, so the chains terminate exactly at lengths three
and five. The rational probabilities obtained from
Eq.~\eqref{eq:spectral-formula} are listed in
Appendix~\ref{app:certificates}; their weighted sums prove the finite-time
claim.

The leading coefficients give a separate short-time certificate:
\begin{equation}\label{eq:right-curvatures}
 \CM(t)=\frac{225}{169}t^2+O(t^4),\qquad
 \KI(t)=\frac{221}{169}t^2+O(t^4).
\end{equation}
They are $\norm{[H_R,\rho_R]}_{\HS}^2/P_R$ and
$2\Var_{\rho_R}(H_R)$, respectively. A norm estimate for the Taylor
remainder, proved in Appendix~\ref{app:technical}, gives
\begin{equation}\label{eq:right-interval}
 \CM(t)-\KI(t)\ge\frac4{169}t^2-64t^3>0,
 \qquad 0<t<\frac1{2704}.
\end{equation}
\end{proof}

These are separate witnesses for separate inequalities. No assertion is
made that one displayed state violates both sides, or that every mixed state
violates either side.

\subsection{Qubits and minimal physical dimension}
\begin{theorem}[Qubit hierarchy]\label{thm:qubit}
For every qubit density matrix, every time-independent Hermitian Hamiltonian,
and every real time, $\SU(t)\le\CM(t)\le\KI(t)$. Consequently, physical
dimension three is minimal for failure of either side of
Eq.~\eqref{eq:source-hierarchy}.
\end{theorem}
\begin{proof}
Let $p_1,p_2$ be the density-matrix eigenvalues, and let $\theta$ be the
polar rotation angle between its eigenbasis and the energy basis. Put
$z=(p_1-p_2)^2$ and $c=\cos^2\theta$, both in $[0,1]$. With
$\tau=(E_2-E_1)t$, the source's qubit formulas
\cite[Eq.~(18), Supplemental Eqs.~(S.24)--(S.31)]{DasMori2026} have the
three-atom form
\begin{equation}\label{eq:qubit-mus}
 \begin{aligned}
 \SU&=F(\mu_S;\tau),&\mu_S&=\frac{(1-\sqrt{1-z})(1-c)}2,\\
 \CM&=F(\mu_K;\tau),&\mu_K&=\frac{z(1-c)}{1+z},\\
 \KI&=F(\mu_I;\tau),&\mu_I&=\frac{1-zc}{2}.
 \end{aligned}
\end{equation}
All three parameters lie in $[0,1/2]$. The lower ordering follows, for
nonzero $\mu_S$, from
\begin{equation}\label{eq:qubit-order}
 \frac{\mu_K}{\mu_S}
 =\frac{(\sqrt{p_1}+\sqrt{p_2})^2}{p_1^2+p_2^2}\ge1,
 \qquad
 \mu_I-\mu_K=\frac{(1-z)(1+zc)}{2(1+z)}\ge0.
\end{equation}
When $\mu_S=0$, the first ordering follows directly from
Eq.~\eqref{eq:qubit-mus}. The function $F(\mu;\tau)$ is nondecreasing in
$\mu$ on $[0,1/2]$, since its derivative is
$\sin^2\tau+8(1-2\mu)\sin^4(\tau/2)\ge0$.
This proves the hierarchy. Pure, maximally mixed, zero-gap, and stationary
endpoints are included by the displayed expressions with their actual
terminated chains. In particular, stationary mixed dynamics need not imply
stationary $I$-purification dynamics, but the ordering still holds.
Dimension one is stationary, and Propositions~\ref{prop:left} and
\ref{prop:right} supply the dimension-three failures.
\end{proof}

\subsection{Failure without a reducing spectator sector}
\begin{proposition}[Upper-side perturbation]\label{prop:robust}
Keep $\rho_R$ fixed and replace $H_R$ by
\begin{equation}\label{eq:perturbed-H}
 H_\delta=\begin{pmatrix}0&1&\delta\\1&0&2\delta\\\delta&2\delta&0\end{pmatrix}.
\end{equation}
For $0<|\delta|<1/\sqrt{135}$, the pair has no nontrivial common reducing
projection and the upper inequality fails at sufficiently small positive
times. At $\delta=1/20$ and $t=1/20000$ one has the explicit certificate
\begin{equation}\label{eq:robust-gap}
 \CM(t)-\KI(t)\ge\frac{378247}{21125000000000000}>0.
\end{equation}
\end{proposition}
\begin{proof}
The exact difference of the quadratic coefficients is
\begin{equation}\label{eq:robust-curvature}
 k_K-k_I=\frac{4-540\delta^2}{169}>0.
\end{equation}
Finite-dimensional analyticity proves the short-time statement. Since
$\rho_R$ has simple spectrum, a common reducing projection must select a
subset of its eigenvectors. Every off-diagonal entry of $H_\delta$ is
nonzero, so no proper nonempty subset reduces it. The finite-time estimate
follows from the explicit norm bound in Appendix~\ref{app:technical}.
\end{proof}

\section{The pure-state factor-two theorem}\label{sec:factor}
The two-sided mixed-state comparison and the pure-density comparison use
different pairs of spectral measures. The latter has a tensor structure
that protects a universal inequality.

\begin{theorem}[Finite-dimensional factor-two inequality]\label{thm:factor}
Let $G$ be a Hermitian matrix and $\psi$ a unit vector in a finite-dimensional
Hilbert space. Set $\psi(t)=e^{-itG}\psi$ and
$\sigma(t)=\ket{\psi(t)}\bra{\psi(t)}$. With the state Krylov chain
generated by $(G,\psi)$ and the operator chain generated by
$([G,\cdot],\sigma(0))$, for every real $t$,
\begin{equation}\label{eq:factor-two}
 \CK(\sigma(t))\ge2\CS(\psi(t)).
\end{equation}
The coefficient two is sharp. Stationary seeds, repeated eigenvalues, and
terminating chains are included.
\end{theorem}
\begin{proof}
Under row vectorization,
\begin{equation}\label{eq:factor-vectorization}
 \chi=\psi\otimes\bar\psi,\qquad
 L=G\otimes I-I\otimes\bar G,\qquad
 \chi(t)=\psi(t)\otimes\overline{\psi(t)}.
\end{equation}
Let $K_i$ be the state Krylov basis and define the polynomial subspaces
\begin{equation}\label{eq:flags}
 V_n=\spanof\{L^k\chi:0\le k\le n\},\qquad
 W_n=\spanof\{K_i\otimes\bar K_j:i+j\le n\}.
\end{equation}
The binomial expansion of $L^k$ proves $V_n\subseteq W_n$. This inclusion
remains valid if a chain terminates, or if the joint cyclic subspace is a
proper subspace of the tensor-product space. No completion of the joint
Krylov basis is needed.

Write $p_i=|\bra{K_i}\psi(t)\rangle|^2$ for the state probabilities and
$p_k^{\rm op}$ for the operator probabilities. Conjugate evolution has the
same probabilities; the possible signs $(-1)^j$ of its Krylov vectors do
not affect their spans. Nested orthogonal projections give
\begin{equation}\label{eq:tail-domination}
 \sum_{k>n}p_k^{\rm op}
 =1-\norm{P_{V_n}\chi(t)}^2
 \ge1-\norm{P_{W_n}\chi(t)}^2
 =\sum_{i+j>n}p_i p_j.
\end{equation}
Summing these nonnegative tails yields
\begin{equation}\label{eq:tail-sum}
 \CK(\sigma(t))\ge\sum_{i,j}(i+j)p_i p_j
 =2\sum_i i p_i=2\CS(\psi(t)).
\end{equation}
For any nonstationary finite-dimensional seed, with
$v=\bra\psi(G-\langle G\rangle)^2\ket\psi>0$, the short-time
coefficients of $\CK$ and $\CS$ are $2v$ and $v$. Their ratio tends to two,
so no larger universal multiplier is possible. A stationary seed gives
zero on both sides.
\end{proof}

The proof also identifies the excess complexity exactly:
\begin{equation}\label{eq:flag-gap}
 \CK(\sigma(t))-2\CS(\psi(t))
 =\sum_{n\ge0}\norm{(P_{W_n}-P_{V_n})\chi(t)}^2.
\end{equation}
In finite dimension the sum terminates. Its positivity is an assertion
about the evolved cyclic vector and the nested subspaces. It does not
require an unqualified positivity claim for a joint number operator on an
unused ambient complement.

\begin{corollary}[The two purification branches]\label{cor:branches}
For either of the source's unitary purification schemes on a fixed finite
enlarged Hilbert space, using its actual normalized initial purification,
\begin{equation}\label{eq:branches}
 \KI(t)\ge2\SI(t),\qquad \KU(t)\ge2\SU(t)
 \qquad(t\in\mathbb R).
\end{equation}
\end{corollary}
\begin{proof}
Apply Theorem~\ref{thm:factor} first with $G=G_I$, and separately with
$G=G_{U^*}$. Both use the pure seed $s$ and the rank-one operator seed
$ss^\dagger$. The two generators need not have the same Krylov structure.
The theorem itself does not require full rank of the original density matrix.
\end{proof}

Equation~\eqref{eq:factor-two} is a special case of the prior
tensor-product superadditivity theorem of Murugan and van Zyl
\cite{MuruganVanZyl2026}: use generators $G$ and $-\bar G$ and the product
seed $\psi\otimes\bar\psi$. Their two subsystem complexities are equal.
The self-contained subspace proof above fixes the source matching without
making the factor-two theorem a novelty claim. We restrict the theorem here
to finite dimension; no infinite-dimensional extension is required for any
counterexample or no-go result in this paper.

\section{Limits of temporal concentration and purity-only control}\label{sec:ratios}
\subsection{The source statement and the quantified questions}
For nonstationary dynamics define, wherever the denominators are nonzero,
\begin{equation}\label{eq:ratios}
 r(t)=\frac{\SU(t)}{\CM(t)},\qquad
 q(t)=\frac{\CM(t)}{\SU(t)}=\frac1{r(t)},\qquad R(t)=Pq(t).
\end{equation}
The main-text Conjecture~2 of Ref.~\cite{DasMori2026} concerns $r$.
The supplement also considers the reciprocal $q$, and its Eq.~(S.32)
defines the purity-rescaled reciprocal $R$. These quantities must not be
interchanged when formulating a concentration statement.

For a positive time series with common recurrence period $T$, use
\begin{equation}\label{eq:cv-definition}
 \avg f=\frac1T\int_0^T f(t)\,dt,\qquad
 \CV(f)^2=\frac{\avg{f^2}}{\bigl(\avg f\bigr)^2}-1.
\end{equation}
The source uses the temporal standard deviation divided by the mean over
a recurrence period. Its roughly five-percent observation refers to a
particular Werner-state scan, not a uniform theorem. Its tabulated relative
maximum fluctuation is a different quantity from the coefficient of
variation. Nor does the source state an exact identity $r=f(P)$.

We consider three explicitly stronger universal formulations. First, at
fixed dimension and fixed Hamiltonian, could one finite constant bound
$\CV(r)$ for every admissible nonstationary full-rank state? Second, could
such a constant bound $\CV(q)$, and hence $\CV(R)$? Third, at fixed
dimension, purity, and Hamiltonian, could $r$ or its period mean have a
finite upper bound determined by purity alone? A separate construction
answers each question negatively. Thus the results are no-go theorems for
these source-motivated universal formulations. They do not assign a truth
value to every possible reading of the qualitative conjecture.

\subsection{Main-ratio concentration}
\begin{theorem}[Unbounded recurrence-period variation of $r$]\label{thm:main-ratio}
Fix the four-dimensional Hamiltonian and two-level block
\begin{equation}\label{eq:main-family}
 H=\diag(0,2,3,4),\qquad
 B=\frac1{10}\begin{pmatrix}5&3\\3&5\end{pmatrix},\qquad
 \rho_\eps=(1-\eps)B\oplus\eps B,
 \quad 0<\eps\le\frac14.
\end{equation}
These states are full rank. With the uniform time average over $T=2\pi$,
and with removable continuation at common recurrences,
\begin{equation}\label{eq:main-cv}
 \CV(r)^2\ge\frac{289}{132710400\pi\eps}-1
 \ \longrightarrow\ \infty\qquad(\eps\to0^+).
\end{equation}
In particular, the rational choice $\eps=10^{-10}$ gives $\CV(r)>80$.
\end{theorem}
\begin{proof}
The eigenvalues are $4(1-\eps)/5,(1-\eps)/5,4\eps/5,\eps/5$.
Let $D_0=(1-\eps)^2+\eps^2$, so $P=17D_0/25$, and note that
$\sqrt B=\left(\begin{smallmatrix}3&1\\1&3\end{smallmatrix}\right)/(2\sqrt5)$.
The normalized return losses are
\begin{equation}\label{eq:main-losses}
 \begin{aligned}
 D_S&=1-\Tr\bigl(\sqrt{\rho(t)}\sqrt\rho\bigr)
 =\frac{(1-\eps)(1-\cos2t)+\eps(1-\cos t)}{10},\\
 D_K&=1-\frac{\Tr(\rho(t)\rho)}P
 =\frac9{34D_0}\bigl[(1-\eps)^2(1-\cos2t)+\eps^2(1-\cos t)\bigr].
 \end{aligned}
\end{equation}
Both chains have exactly the five frequencies $0,\pm1,\pm2$ with positive
weights. Lemma~\ref{lem:return} yields $f_\eps/8\le r\le8f_\eps$,
where, writing $c=\cos^2(t/2)$,
\begin{equation}\label{eq:main-f}
 \begin{aligned}
 f_\eps=\frac{D_S}{D_K}
 &=\frac{17D_0}{45}
 \frac{4(1-\eps)c+\eps}{4(1-\eps)^2c+\eps^2}\\
 &=\frac{17D_0}{45(1-\eps)}
 \left[1+\frac{\eps(1-2\eps)}{4(1-\eps)^2c+\eps^2}\right].
 \end{aligned}
\end{equation}
The common recurrence factor has been cancelled. The elementary integral
in Appendix~\ref{app:technical} gives
\begin{equation}\label{eq:main-mean}
 \avg{f_\eps}=\frac{17D_0}{45(1-\eps)}
 \left[1+\frac{1-2\eps}{\sqrt{4(1-\eps)^2+\eps^2}}\right]
 \le\frac{68}{81}<1,\qquad \avg r\le8.
\end{equation}
On the window $|t-\pi|\le\eps$, one has $c\le\eps^2/4$ and
$D_0\ge1/2$. Consequently
\begin{equation}\label{eq:main-window}
 f_\eps\ge\frac{17}{180\eps},\qquad
 r\ge\frac{17}{1440\eps},\qquad
 \avg{r^2}\ge\frac{289}{2073600\pi\eps}.
\end{equation}
Combining the mean and second-moment bounds gives
Eq.~\eqref{eq:main-cv}. At $\eps=10^{-10}$, the elementary estimate
$\pi<22/7$ makes its lower bound exceed $80^2$.
\end{proof}

At $t=\pi$ the heavily populated block has returned, while the light block
has not. The light coherent sector contributes order $\eps$ to the
$\sqrt\rho$ measure and order $\eps^2$ to the normalized $\rho$ measure.
The resulting excursion is narrow, but its second moment is large enough
to rule out any finite uniform coefficient of variation.

\subsection{Reciprocal and rescaled-reciprocal concentration}
Inversion does not preserve relative temporal concentration. The preceding
family therefore cannot substitute for a proof about the supplement's
ratio.

\begin{theorem}[Unbounded recurrence-period variation of $q$]\label{thm:reciprocal}
Keep $H=\diag(0,2,3,4)$. Let $\eta>0$, put $\eps=\eta^2$, and assume
$0<\eps\le1/16$. Define
\begin{equation}\label{eq:reciprocal-family}
 A_\eta=\eta\begin{pmatrix}2&1\\1&2\end{pmatrix}
 \oplus\begin{pmatrix}1&\eta^3\\\eta^3&1\end{pmatrix},
 \qquad \rho_\eta=\frac{A_\eta^2}{N},
 \qquad N=\Tr A_\eta^2.
\end{equation}
For these full-rank states, using $T=2\pi$ and removable recurrence values,
\begin{equation}\label{eq:reciprocal-cv}
 \CV(q)^2\ge\frac1{1679616\pi\eps}-1
 \ \longrightarrow\ \infty,\qquad \CV(R)=\CV(q).
\end{equation}
The rational choice $\eta=1/100000$ gives $\CV(q)>40$.
\end{theorem}
\begin{proof}
The eigenvalues of $A_\eta$ are $3\eta,\eta,1+\eta^3,1-\eta^3$,
which are positive in the stated range. Its normalization and fourth trace
are
\begin{equation}\label{eq:reciprocal-traces}
 N=2+10\eps+2\eps^3,\qquad
 T_4=\Tr A_\eta^4=2+82\eps^2+12\eps^3+2\eps^6,
 \qquad P=\frac{T_4}{N^2}.
\end{equation}
The positive square root of $\rho_\eta$ is $A_\eta/\sqrt N$. Write
$a=1-\cos2t$ and $b=1-\cos t$. Direct autocorrelation gives
\begin{equation}\label{eq:reciprocal-losses}
 D_S=\frac{2\eps}{N}a+\frac{2\eps^3}{N}b,\qquad
 D_K=\frac{32\eps^2}{T_4}a+\frac{8\eps^3}{T_4}b.
\end{equation}
Again both chains have five positive-weight atoms. Thus
$g_\eps/8\le q\le8g_\eps$, with
\begin{equation}\label{eq:reciprocal-g}
 g_\eps=\frac{D_K}{D_S}
 =\frac{4N\eps}{T_4}\frac{16c+\eps}{4c+\eps^2},
 \qquad c=\cos^2(t/2).
\end{equation}
On the specified range, $2\le N,T_4\le3$, and the period integral yields
\begin{equation}\label{eq:reciprocal-mean}
 \avg{g_\eps}=\frac{4N\eps}{T_4}
 \left[4+\frac{1-4\eps}{\sqrt{4+\eps^2}}\right]
 \le27\eps,\qquad \avg q\le216\eps.
\end{equation}
On $|t-\pi|\le\eps$, where $c\le\eps^2/4$,
\begin{equation}\label{eq:reciprocal-window}
 g_\eps\ge\frac{2N}{T_4}\ge\frac43,\qquad
 q\ge\frac16,\qquad \avg{q^2}\ge\frac{\eps}{36\pi}.
\end{equation}
These inequalities prove Eq.~\eqref{eq:reciprocal-cv}; $\pi<22/7$ verifies
the finite rational witness. Finally $P$ is positive and independent of
time, so multiplication by $P$ leaves the coefficient of variation unchanged.
\end{proof}

Here the mean becomes small while finite excursions persist in a narrowing
window. A decreasing absolute fluctuation is therefore compatible with
unbounded fluctuation relative to the mean.

\subsection{Recurrences and admissibility}
In both temporal families, the return losses are positive combinations of
$1-\cos t$ and $1-\cos2t$. They vanish together only at
$t\in2\pi\mathbb Z$. For positive-matrix returns, the lower bound in
Lemma~\ref{lem:return} then implies that these are precisely the
zero-complexity points. At each such point the two complexities have
positive quadratic coefficients, so the ratio has a finite positive
removable limit. The point $t=\pi$ is not a global recurrence: the gap-one
block remains active for every allowed positive parameter. Each temporal
integral is finite at each fixed parameter. The divergence conclusions use
the displayed inequalities and do not exchange a singular state limit
with the time integral.

\subsection{Fixed purity}
\begin{theorem}[No finite purity-only upper control]\label{thm:purity}
Fix $d=4$, $H=\diag(2,3,0,1)$, and $P=1/2$. For $0<\eps<5/18$, let
\begin{equation}\label{eq:purity-family}
 a_\pm=\frac{1-\eps\pm\sqrt{2\eps-59\eps^2/25}}2,
 \qquad \rho_\eps=\diag(a_+,a_-)\oplus\eps B,
\end{equation}
with $B$ as in Eq.~\eqref{eq:main-family}. These states are full rank and
have exactly the specified purity. Their main ratio satisfies, at every
real time with removable continuation at recurrences,
\begin{equation}\label{eq:purity-bounds}
 \frac5{18\eps}\left(1-\frac\eps{10}\right)
 \le r_\eps(t)\le\frac5{18\eps}.
\end{equation}
Thus $r_\eps$ and its recurrence-period mean diverge uniformly as
$\eps\to0^+$, whereas $q_\eps$ and $Pq_\eps$ tend uniformly to zero.
\end{theorem}
\begin{proof}
The eigenvalues are $a_+,a_-,4\eps/5,\eps/5$. The radicand is positive,
and
\begin{equation}\label{eq:purity-positive}
 (1-\eps)^2-(2\eps-59\eps^2/25)
 =(1-14\eps/5)(1-6\eps/5)>0.
\end{equation}
Therefore $a_->0$. Their sum is one and their squared sum is $1/2$.
Only the last block carries coherence, with gap one. Its total
nonzero-gap weights and exact ratio are
\begin{equation}\label{eq:purity-ratio}
 \begin{aligned}
 \mu_K&=\frac{9\eps^2}{25},\qquad \mu_S=\frac\eps{10},
 \qquad x=\sin^2(t/2),\\
 r_\eps(t)&=\frac{F(\mu_S;t)}{F(\mu_K;t)}
 =\frac5{18\eps}
 \frac{1+(1-\eps/5)x}{1+(1-18\eps^2/25)x}.
 \end{aligned}
\end{equation}
The reduced denominator is positive. For the stated parameter interval
the last fraction decreases with $x\in[0,1]$; its endpoint estimates
give Eq.~\eqref{eq:purity-bounds}. At $t\in2\pi\mathbb Z$ the removable
value is $5/(18\eps)$. The uniform bounds also apply to the period mean.
\end{proof}

Because dimension, purity, and Hamiltonian are all fixed, the family
excludes even a finite upper bound allowed to depend on those fixed data.
In particular it excludes finite purity-only upper control of $r$ or its
mean and positive purity-only lower control of $q$ or $Pq$.
It leaves intact the weaker statement that purity influences these
quantities together with additional spectral information. It does not
contradict the reported behavior within the particular Werner family.

\section{Restricted survivors and the spectral mechanism}\label{sec:survivors}
\subsection{A uniform qubit bound for relative variation}
The failure of uniform concentration in dimension four coexists with a
restricted theorem for qubits.

\begin{proposition}[Qubit coefficient-of-variation bound]\label{prop:qubit-cv}
For a nonstationary qubit under a time-independent Hermitian Hamiltonian,
extend $r,q,R$ through their removable recurrence values. For any
probability-weighted time average,
\begin{equation}\label{eq:qubit-cv}
 \CV(r),\ \CV(q),\ \CV(R)
 \ \le\ \frac{\kappa_*-1}{2\sqrt{\kappa_*}},
 \qquad \kappa_*=\frac{17+7\sqrt7}{27}.
\end{equation}
The upper bound is approximately $0.1375633885$. No claim is made that
this value is attained by a qubit trajectory and averaging measure.
\end{proposition}
\begin{proof}
Put $y=2\sqrt{p_1p_2}\in[0,1)$ and $h=\sin^2\theta>0$; nonstationarity
also requires a nonzero energy gap. The coefficients in
Eq.~\eqref{eq:qubit-mus} become
\begin{equation}\label{eq:qubit-range}
 \mu_S=\frac{(1-y)h}{2},\qquad
 \mu_K=\frac{(1-y^2)h}{2-y^2},\qquad
 \kappa=\frac{\max_t r(t)}{\min_t r(t)}
 =\frac{1-\mu_S}{1-\mu_K}.
\end{equation}
The same range factor applies to $q$ and $R$. It increases with $h$, so
\begin{equation}\label{eq:kappa-max}
 \kappa\le\frac{(1+y)(2-y^2)}2
 \le\frac{17+7\sqrt7}{27},\qquad
 y_{\max}=\frac{\sqrt7-1}{3}.
\end{equation}
The second derivative of the middle expression is $-1-3y<0$, and its
endpoint values are one. For a positive random variable $X\in[m,M]$,
the inequality $\mathbb E[(X-m)(M-X)]\ge0$ implies
\begin{equation}\label{eq:bounded-variance}
 \Var(X)\le(M-\mathbb EX)(\mathbb EX-m),\qquad
 \CV(X)\le\frac{M-m}{2\sqrt{mM}}
 =\frac{\kappa-1}{2\sqrt\kappa}.
\end{equation}
The last step maximizes the relative variance over the allowed mean.
Apply this bound to each of the three ratios.
\end{proof}

Completely stationary $0/0$ ratios are excluded from this proposition.
Dimension four suffices for the temporal no-go constructions, while the
qubit bound excludes dimension two. Whether dimension three suffices for
unbounded temporal variation is not settled here. No minimum-dimension
claim is made for the fixed-purity no-go theorem.

\subsection{State-dependent comparisons}
The return bound yields useful comparisons retaining the two spectral
measures. Let $m_S,m_K$ be the finite lengths of the $\SU$ and $\CM$
chains and define
\begin{equation}\label{eq:returns}
 A_S=\Tr\bigl(\sqrt{\rho(t)}\sqrt\rho\bigr),\qquad
 A_K=\frac{\Tr(\rho(t)\rho)}P.
\end{equation}
\begin{corollary}[Return comparison]\label{cor:return-comparison}
When both return losses are nonzero,
\begin{equation}\label{eq:state-dependent}
 \frac{1-A_S^2}{(m_K-1)(1-A_K^2)}
 \le r(t)\le
 \frac{(m_S-1)(1-A_S^2)}{1-A_K^2}.
\end{equation}
\end{corollary}
\begin{proof}
Divide the lower return bound for the numerator by the upper one for
the denominator, and conversely.
\end{proof}

These bounds involve return amplitudes and cyclic dimensions, rather
than purity alone. Similarly, a necessary short-time condition for the
source hierarchy would be
\begin{equation}\label{eq:curvature-necessary}
 \norm{[H,\sqrt\rho]}_{\HS}^2
 \le\frac{\norm{[H,\rho]}_{\HS}^2}{P}
 \le2\Var_\rho(H),\qquad
 \Var_\rho(H)=\Tr(\rho H^2)-(\Tr\rho H)^2.
\end{equation}
The displayed counterexamples violate the respective comparisons.
Satisfying these necessary curvature inequalities is not claimed to be
sufficient for an all-time hierarchy.

\subsection{Why the comparisons separate}
The mixed seed and the square-root seed give different powers of a
population weight to a coherent sector. Hilbert--Schmidt normalization
then introduces the global purity into only the former. A highly populated
stationary sector can suppress mixed-state operator complexity more strongly
than purification spread complexity. Independent block frequencies permit
partial recurrences that amplify this mismatch in time. Holding the total
purity fixed does not fix the distribution of population among these sectors.

By comparison, a pure density matrix is a tensor product after
vectorization. Its difference-convolution measure comes with the nested
polynomial subspaces of Section~\ref{sec:factor}. No corresponding
containment is generally available for the $\rho$ and $\sqrt\rho$
measures. Qubits are protected by their common three-atom form and the
monotone ordering of the three coefficients. The loss of these structures,
rather than an inconsistency between state and operator definitions,
explains the different outcomes.

\section{Discussion}\label{sec:discussion}
Purification preserves the reduced density matrix but can change the
spectral measure defining complexity. With the source's normalized seeds,
each side of the mixed/purification hierarchy fails in an exact full-rank
qutrit example. The qubit theorem fixes the minimum physical dimension
for both failures. The upper-side perturbation further shows that its
failure does not require an exactly decoupled spectator.

The ratio results delimit a different issue. Uniform temporal
concentration of $r$ and of its reciprocal fail by different four-level
families, and fixed purity cannot supply finite universal upper control
of $r$ or its mean. These statements concern explicit universal
formalizations of the qualitative source discussion. They neither replace
that discussion by an invented quantitative conjecture nor invalidate its
sample-specific numerical observations. Typical-state statements,
restricted spectral families, and assumptions controlling the smallest
populations would require their own quantified formulations.

The surviving factor-two inequality and qubit bounds demonstrate that
useful comparisons remain possible when their spectral structure is
specified. The pure-density inequality belongs to prior tensor-product
superadditivity and has been included to separate that protected
comparison from the failed mixed-state ordering. The state-dependent
return bounds identify information that a replacement comparison can
retain. No universal classification of such replacements is asserted.

All decisive constructions are finite dimensional. The temporal
minimum-dimension question between three and four remains open within
this analysis, and no claim is made about unrestricted infinite-dimensional
or nonunitary dynamics. The accompanying exact certificates make the
stated finite results directly checkable without treating numerical plots
as proofs.

\appendix
\section{Exact finite certificates}\label{app:certificates}
\subsection{Complete probabilities for the two witnesses}
For the lower-side witness at $t=\pi$, the probability vectors are
\begin{equation}\label{eq:left-probabilities}
 p^{\mathrm M}=\left(\frac{7744}{8281},0,\frac{537}{8281}\right),
 \qquad
 p^{S,U^*}=\left(\frac{400}{441},0,\frac{41}{441}\right).
\end{equation}
Each has unit sum. Multiplying the final entries by two gives the
complexities in Eq.~\eqref{eq:left-values}.

For the upper-side witness, the squared nonzero Lanczos coefficients are
\begin{equation}\label{eq:right-lanczos}
 (b_n^2)_{\mathrm M}=\left(\frac{225}{169},\frac{451}{169}\right),
 \qquad
 (b_n^2)_{K,I}=\left(\frac{17}{13},\frac{43}{26},
                      \frac{1377}{1118},\frac{451}{559}\right).
\end{equation}
All diagonal coefficients vanish. At $t=\pi/3$ the full probabilities are
\begin{equation}\label{eq:right-probabilities}
 \begin{aligned}
 p^{\mathrm M}&=\left(\frac{458329}{1827904},\frac{675}{2704},
                           \frac{913275}{1827904}\right),\\
 p^{K,I}&=\left(\frac{1500625}{7311616},\frac{62475}{140608},
 \frac{45778977}{157199744},\frac{7803}{140608},\frac{1173051}{314399488}\right).
 \end{aligned}
\end{equation}
For a direct finite check, each row sums to one and its index-weighted
sum is the corresponding rational number in Eq.~\eqref{eq:right-values}.
These data follow from the exact measures in Eq.~\eqref{eq:right-measures}.
The terminal polynomials are $x(x-2)(x+2)$ and
$x(x-2)(x-1)(x+1)(x+2)$, respectively, so no floating-point stopping
criterion is involved. The authoritative certificate additionally compares
these spectral calculations with direct matrix-commutator Lanczos
calculations of every probability.

\subsection{An unbounded qutrit ratio family}\label{app:unbounded-left}
\begin{proposition}[No fixed positive multiplier for the lower side]\label{prop:unbounded-left}
For every real $m\ge4$, keep $H_L=\diag(2,0,1)$ and take
\begin{equation}\label{eq:m-family}
 \rho_m=\frac1{m^2+5}
 \begin{pmatrix}m^2&0&0\\0&5/2&3/2\\0&3/2&5/2\end{pmatrix}.
\end{equation}
The state is full rank, and $\SU(t)>\CM(t)$ for every
$t\notin2\pi\mathbb Z$. Moreover,
\begin{equation}\label{eq:m-asymptotic}
 \frac{\SU(\pi)}{\CM(\pi)}\sim\frac{m^2}{9}
 \longrightarrow\infty\qquad(m\to\infty).
\end{equation}
Thus no constant $c>0$ can make $c\SU\le\CM$ hold for all full-rank
qutrit states, Hermitian Hamiltonians, and real times.
\end{proposition}
\begin{proof}
The spectrum is $(m^2,4,1)/(m^2+5)$ and
\begin{equation}\label{eq:m-weights}
 P_m=\frac{m^4+17}{(m^2+5)^2},\qquad
 \mu_S=\frac1{2(m^2+5)},\qquad
 \mu_K=\frac9{2(m^4+17)}.
\end{equation}
Their difference has numerator $m^4-9m^2-28$, which is positive for
$m\ge4$. Both complexities are given by $F$, and
\begin{equation}\label{eq:F-difference}
 F(a;t)-F(b;t)
 =(a-b)\bigl[\sin^2t+8(1-a-b)\sin^4(t/2)\bigr].
\end{equation}
Here $a=\mu_S>b=\mu_K$ and $a+b<1$, proving the all-time strict
inequality away from recurrences. At $\pi$, $F(\mu;\pi)=8\mu(1-\mu)$,
which gives Eq.~\eqref{eq:m-asymptotic}. Every finite $m$ is an admissible
full-rank state; the limiting rank-deficient state is not used as a witness.
\end{proof}

\subsection{A rational fixed-purity anchor}
At $\eps=1/10$, the family in Theorem~\ref{thm:purity} is
\begin{equation}\label{eq:purity-anchor}
 \rho_{1/10}=\frac1{100}
 \begin{pmatrix}66&0&0&0\\0&24&0&0\\0&0&5&3\\0&0&3&5\end{pmatrix},
 \qquad P=\frac12.
\end{equation}
This supplies a finite rational comparison state at the same dimension,
purity, and Hamiltonian as the entire family. It prevents confusing the
fixed-purity conclusion with an effect obtained solely by changing purity.

\section{Technical bounds}\label{app:technical}
\subsection{Controlled finite-time Taylor remainders}
For a complexity written in its finite cyclic space as
$C(t)=\bra v e^{itL}N e^{-itL}\ket v$, with $N\ket v=0$, repeated
commutators give
\begin{equation}\label{eq:taylor-norm}
 |C'''(t)|\le(2\norm L)^3\norm N,\qquad
 |C(t)-kt^2|\le\frac{(2\norm L)^3\norm N}{6}|t|^3,
\end{equation}
where $k$ is the quadratic coefficient. For Eq.~\eqref{eq:right-state},
both Liouvillian norms are at most two and the number-operator norms are
two and four. Adding the two remainder estimates yields $64|t|^3$,
proving Eq.~\eqref{eq:right-interval}.

For $H_\delta$ at $\delta=1/20$, the maximal absolute row sum bounds
$\norm{H_\delta}\le11/10$. Each operator chain has at most seven
distinct energy-gap atoms, so each number-operator norm is at most six.
The two Liouvillian norms are at most $11/5$. Hence
\begin{equation}\label{eq:perturbed-remainder}
 \CM(t)-\KI(t)
 \ge\frac{53}{3380}t^2-\frac{21296}{125}|t|^3.
\end{equation}
At $t=1/20000$ this is exactly the positive rational lower bound in
Eq.~\eqref{eq:robust-gap}. These are conservative finite-time bounds,
not numerical estimates of the exact differences.

\subsection{Period integral and mean estimates}
For $a,b>0$, tangent substitution after splitting the interval at $\pi$
gives
\begin{equation}\label{eq:integral}
 \frac1{2\pi}\int_0^{2\pi}
 \frac{dt}{a\cos^2(t/2)+b}
 =\frac1{\sqrt{b(a+b)}}.
\end{equation}
For the main-ratio family, $D_0\le1$, $1-\eps\ge3/4$, and
$\sqrt{4(1-\eps)^2+\eps^2}\ge2(1-\eps)$ bound the mean in
Eq.~\eqref{eq:main-mean} by $68/81$. On the peak window, its rational
denominator is at most $2\eps^2$ while its numerator is at least
$\eps$; together with $D_0\ge1/2$ this proves
Eq.~\eqref{eq:main-window}.

For the reciprocal family, the positive coefficients in
Eq.~\eqref{eq:reciprocal-traces} show that $2\le N,T_4\le3$ on
$0<\eps\le1/16$. The bracket in Eq.~\eqref{eq:reciprocal-mean} is at
most $9/2$, so its prefactor bounds the mean by $27\eps$.
On the window the denominator $4c+\eps^2$ is at most $2\eps^2$
and the numerator $16c+\eps$ is at least $\eps$. This gives
$g_\eps\ge2N/T_4\ge4/3$ and hence the second-moment estimate.

\subsection{The fixed-purity endpoint bound}
In Eq.~\eqref{eq:purity-ratio} the fractional-linear factor decreases
because $\eps/5>18\eps^2/25$ for $0<\eps<5/18$. Its value at
$x=0$ is one. At $x=1$ it is
\begin{equation}\label{eq:purity-endpoint}
 \frac{1-\eps/10}{1-9\eps^2/25}\ge1-\eps/10.
\end{equation}
This gives the lower bound in Eq.~\eqref{eq:purity-bounds}.

\section{Certificate map and reproducibility}\label{app:reproducibility}
The preserved internal manuscript archive contains the complete supplied input bundle,
integrity records, and certification materials. The public repository provides a curated
reproducibility view and does not redistribute all internal archival or third-party materials. The scientific
authority is the archive
\path{Krylov_PRL_WORK_FREEZE_2026-09-23_v3.zip}.
Its SHA-256 is recorded in the project's integrity report and checked
against the externally specified anchor before use. The nested second-run
archive and first-run certificate retain their original bytes.

For a compact map, the first-run certificate is the V3 directory
\path{prior_left/P2_3_Krylov_certificate/}; the second-run certificate
directory is
\path{current_final/Krylov_remaining_attack/certificates/}.
The former supplies the lower-side qutrit and its unbounded family.
Within the latter, \path{RIGHT_BOUND_ONE_PAGE.md} and
\path{RIGHT_BOUND_DETAILS.md} supply the upper-side witness,
its perturbation, and minimality. \path{FACTOR_TWO_THEOREM.md}
supplies the polynomial-subspace proof. \path{TARGET2_TEMPORAL_NOGO.md}
and \path{TARGET2_PURITY_NOGO.md} supply the ratio families and all
constants retained here. \path{RESTRICTED_SURVIVORS.md} supplies
the qubit variation and return bounds. Exact probability data are retained
in the first-run results and the second-run
\path{results/target1r_exact.json}. The accompanying
\path{CLAIM_SOURCE_MAP.md} gives complete paths and the scope of
each theorem and headline statement.

The supplied certification evidence includes two adversarial AI audits and
two frozen AI clean-room reproductions of seven specified core targets.
Both clean-room runs reproduce the two finite qutrit certificates,
the qubit hierarchy, the finite-dimensional factor-two result, and
the conclusions of the three ratio no-go constructions. Their task was
to reproduce supplied witnesses and families, not independently discover
those inputs. This is AI certification evidence; it is not independent
human validation or peer review. Secondary byproducts are not silently
included in the clean-room coverage.

Different clean-room estimates of concentration constants are recorded
only in \path{PROPOSED_NONAUTHORITATIVE_IMPROVEMENTS.md}. All theorem
constants in this manuscript remain those of V3. A stronger estimate would
require a separate targeted audit and a new scientific freeze before
promotion into the claim set.

The adversarial audits identified a stale-manifest defect in the expanded
tree of the historical V2 packaging. V3 records the packaging repair and
preserves the unchanged authoritative mathematical archives. Replays must
use a disposable extraction: the original combined verifier writes a
results file and therefore changes a frozen tree if run in place.
The manuscript project keeps its input archives immutable. The certification package described here did not itself
perform any public-release, submission, or author-contact
action.

\section*{AI contribution statement}
AI systems played the primary and largest role in the scientific work reported here. They identified the source paper as a research target, developed the attacks on its conjectures, constructed the counterexamples, derived the proofs and auxiliary results, produced and checked the computational certificates, and assisted in the adversarial audits, clean-room reproductions, and preparation of the manuscript. Dehao Lin designed and supervised the overall research workflow, made the project and release decisions, reviewed the resulting artifacts, and assumes responsibility for the final manuscript.

\paragraph{Funding.} This research received no specific funding.
\paragraph{Acknowledgements.} None.
\paragraph{License.} The original manuscript is licensed under
\href{https://creativecommons.org/licenses/by/4.0/}{Creative Commons Attribution 4.0 International (CC BY 4.0)}.
This notice does not relicense third-party material.
\bibliographystyle{plain}
\bibliography{references}
\end{document}
